
% ДЛЯ КА ВВ и КА СВ! КА СВ отличие только во входах на включение выключателя в схеме В используются xlsx VV_SwitchDevice и SV_SwitchDevice
\needspace{3\baselineskip}
\color{uniblue}{\subsection{Сигнализация}}
\color{black}

\begin{enumerate}[label=\arabic{section}.\arabic{subsection}.\arabic{enumi}, labelsep=4pt, leftmargin=0pt, itemindent=57pt, itemsep=0pt, parsep=5pt] % если что можно крутить itemsep


%\begin{enumerate}[label=\arabic{section}.\arabic{subsection}.\arabic{enumi}.\arabic*, labelsep=4pt, leftmargin=0em, itemindent=65pt, parsep=0pt]
        
%%%%%%%%%%%%%%%%%%%%%%%%%%%%%%%%%%%%%%%%%%%%%%%%%%%%%%%%%%%%%%%%%%%%% ОБЩИЕ СВЕДЕНИЯ %%%%%%%%%%%%%%%%%%%%%%%%%%%%%%%%%%%%%%%%%%%%%%%%%%%%%%%%%%%%%%

\item Функция сигнализации (см. рисунок 2.15) принимает сигналы неисправностей, контроля положения испытательных блоков и двери. Функция формирует следующие обобщенные сигналы:
    \begin{itemize}
        \item вызов;
        \item неисправность;
        \item БИ в нерабочем положении.
    \end{itemize}


\item Ключом «Вывод терминала» блокируется формирование выходных сигналов. 

\small
\begin{longtable}{|>{\centering\arraybackslash}m{5.3cm}|>{\centering\arraybackslash}m{3.3cm}|>{\centering\arraybackslash}m{4.2cm}|>{\centering\arraybackslash}m{1.8cm}|>{\centering\arraybackslash}m{1cm}|}
\caption{Параметры для настройки функции сигнализации\hfill\vspace{-0.5\baselineskip}}\label{mtzvn:tbl1}\\ 
\hline
\rowcolor{gray!30}
Параметр (Параметр на ИЧМ) & Условное обозначение на схеме & Значение/ Диапазон & Единица измерения & Шаг \\ 
\hline
\endfirsthead
%\caption{\emph{Продолжение\hfill\vspace{-0.5\baselineskip}}} \\ % Отступ обозначения таблицы от таблицы и выравнивание надписи слева
\caption*{\hspace{3pt}\emph{Продолжение таблицы \ref{mtzvn:tbl1}\hfill\vspace{-0.5\baselineskip}}} \\ % сделано по ГОСТ 2.105 п.6.8.7
\hline
\rowcolor{gray!30}
Параметр (Параметр на ИЧМ) & Условное обозначение на схеме & Значение/ Диапазон & Единица измерения & Шаг \\ 
%\hline
\endhead
%\hline
\endfoot
%\hline
\endlastfoot


%==+t1*PLGSOCC1|_TTCC6> ПДС НКУ
\multicolumn{5}{|c|}{ ПДС НКУ } \\ \hline 
\centering Контроль положения блока испытательного №1 (КонтрSG1) & \centering КонтрSG1 & \centering 0 = Не предусмотрено\\1 = Предусмотрено & \centering -- & \centering \arraybackslash -- \\
\hline
\centering Контроль положения блока испытательного №2 (КонтрSG2) & \centering КонтрSG2 & \centering 0 = Не предусмотрено\\1 = Предусмотрено & \centering -- & \centering \arraybackslash -- \\
\hline
\centering Контроль положения блока испытательного №3 (КонтрSG3) & \centering КонтрSG3 & \centering 0 = Не предусмотрено\\1 = Предусмотрено & \centering -- & \centering \arraybackslash -- \\
\hline
\centering Контроль положения блока испытательного №4 (КонтрSG4) & \centering КонтрSG4 & \centering 0 = Не предусмотрено\\1 = Предусмотрено & \centering -- & \centering \arraybackslash -- \\
\hline
\centering Контроль положения блока испытательного №5 (КонтрSG5) & \centering КонтрSG5 & \centering 0 = Не предусмотрено\\1 = Предусмотрено & \centering -- & \centering \arraybackslash -- \\
\hline
\centering Контроль положения блока испытательного №6 (КонтрSG6) & \centering КонтрSG6 & \centering 0 = Не предусмотрено\\1 = Предусмотрено & \centering -- & \centering \arraybackslash -- \\
\hline
\centering Контроль положения блока испытательного №7 (КонтрSG7) & \centering КонтрSG7 & \centering 0 = Не предусмотрено\\1 = Предусмотрено & \centering -- & \centering \arraybackslash -- \\
\hline
\centering Контроль положения двери (КонтрДвери) & \centering КонтрДвери & \centering 0 = Не предусмотрено\\1 = Предусмотрено & \centering -- & \centering \arraybackslash -- \\
\hline
\centering Контроль оперативного тока технологической сигнализации (КонтрОТ) & \centering КонтрОТ & \centering 0 = Не предусмотрено\\1 = Предусмотрено & \centering -- & \centering \arraybackslash -- \\
\hline
\centering Выдержка времени срабатывания сигнализации (Тсигн) & \centering Тсигн & \centering 1 ... 100000 & \centering мс & \centering \arraybackslash 1 \\
\hline
%===t1




\end{longtable}
\normalsize











% % ФОРМУЛА №1
% \begin{equation*}
% \begin{bmatrix}
% I_{A,n*}\\ 
% I_{B,n*}\\ 
% I_{C,n*}
% \end{bmatrix}= k_{AM,n}\times \begin{bmatrix}
% M_{11,n} & M_{12,n} & M_{13,n}\\ 
% M_{21,n} & M_{22,n} & M_{23,n}\\ 
% M_{31,n} & M_{32,n} & M_{33,n}
% \end{bmatrix}\times \begin{bmatrix}
% I_{A,n}\\ 
% I_{B,n}\\ 
% I_{C,n}
% \end{bmatrix},
% \end{equation*}
% \begin{eqexpl}[25mm]
% \item{$n$}индекс, обозначающий сторону (плечо) ДЗТ;
% \item{$k_{AM,n}$}коэффициент амплитудного выравнивания;
% \item{$I_{A,n},I_{B,n},I_{C,n}$}измеренные фазные токи стороны $n$;
% \item{$M_{11,n}\dots M_{33,n} $}коэффициенты матрицы компенсации фазового сдвига в соответствие  с группой соединения обмоток трансформатора.
% \end{eqexpl}


%%%%%%%%%%%%%%%%%%%%%%%%%%%%%%%%%%%%%%%%%%%%%%%%%%%%%%%%%%% КОНЕЦ Выкатной элемент (ВЭ) %%%%%%%%%%%%%%%%%%%%%%%%%%%%%%%%%%%%%%%%%%%%%%%%%%%%%%%%%%


\end{enumerate}